\part{Diagrammatics}

\label{pt:diagrammatics}

%%%%%%%%%%%%%%%%%%%%%%%%%
\section{Diagrammatics for Standard Bimodules}
\label{sec-diagstd}
%%%%%%%%%%%%%%%%%%%%%%%%%

%===========
\subsection{Diagrammatic definition}
\label{subsec-StdDiagDef}
%===========

We assume that the reader is familiar with diagrammatics for monoidal categories with biadjunctions. An introductions to the topic may be found in \cite{Lauda}.

Fix $(W,S)$. In the paper \cite{EWFenn}, we give a diagrammatic monoidal presentation for the 2-groupoid of $W$ over $\Bbbk$. As we have seen, this is equivalent to the category $\StdBim_0$. We recall that definition, beefing it up to add polynomials and obtain the 2-groupoid of $W$ over $R$, which is equivalent to $\StdBim$. We introduce it in a way which will aid our definition of Soergel diagrammatics in the next chapter.

\begin{defn} \label{defn: graph}(For the purposes of this paper) a \emph{graph with boundary} for $(W,S)$ is an \emph{isotopy class} of a graph with boundary, properly embedded in the planar strip $\R
\times [0,1]$. In other words, all vertices of the graph lie on the interior of the strip, and edges may terminate either at a vertex or on the boundary of the strip. We call the place
where an edge meets the boundary a \emph{boundary point}; boundary points are not vertices. Edges may also form closed loops. We also allow decorations to float in regions cut out by the graph;
these are thought of as 0-valent vertices with labels (we will always clarify in subdefinitions what labels are allowed). Note that a connected component either contains a vertex, is a
loop, is an arc between two boundary points, or is a decoration.
The number of vertices and the number of components are required to be finite. Every edge is ``colored" (i.e. labelled) by
an element of $S$. Isotopy is allowed to move boundary points within the boundary of the strip. The boundary points on $\R \times \{0\}$ (resp. $\{1\}$) give a finite sequence of colored
points, known as the \emph{bottom boundary} (resp. \emph{top}) of the graph. We will often abuse notation and refer to a representative of the isotopy class as a graph. \end{defn}

We fix a realization of $(W,S)$, yielding the polynomial ring $R$.

\begin{defn} A \emph{standard graph} is a graph with boundary for which: the only decorations are boxes labelled by homogeneous $f \in R$; for every vertex there is a pair $s,t \in S$
such that the vertex is $2m_{st}$-valent and its edges alternate in color between $s$ and $t$. The \emph{degree} of the graph is the sum of the degrees of every box. \end{defn}

We draw standard graphs with dashed edges, to distinguish them from the Soergel graphs to be defined later. Here is an example where $S=\{r,b,g\}$ (for \emph{r}ed, \emph{b}lue and \emph{g}reen) and $m_{r,g}=5$, $m_{r,b}=3$, $m_{b,g}=2$:

%psfrag f
\igc{1}{stddiagexample}

\begin{defn} Let $\DStd$ denote the $\Bbbk$-linear additive graded monoidal category defined as follows. Objects are (direct sums of grading shifts of) sequences $\xb$ of elements of
$S$, with monoidal structure given by concatenation. The empty sequence $\emptyset=\1$ is the monoidal identity. We may draw these objects as colored points on a line, assigning one
color to each index in $S$. The space $\Hom_{\DStd}(\xb,\yb)$ will be the free $\Bbbk$-module generated by standard graphs with bottom boundary $\xb$ and top boundary $\yb$, modulo the
relations below. Hom spaces will be graded by the degree of the graphs, and all the relations below are homogeneous.

The first (unwritten) relation states that boxes add and multiply just as homogeneous polynomials in $R$ do. Thus we could have chosen our generating decorations to merely be boxes labelled by generators of $R$. The other tacit relation is that cups and caps form a collection of biadjoint pairs, and that all maps are cyclic with respect to these
biadjunctions, having the appropriate rotational symmetry. This is implied by the definition of a graph, because an isotopy class of diagram unambiguously represents a morphism:

\igc{1}{biadjoint}

The following relation describes how to slide boxes through edges.

%psfrag f soff
\begin{equation} \label{slidepolythrustd} \ig{1}{slidepolythrustd} \end{equation}

The following relations hold for any $s \in S$. As usual, empty space here denotes the identity endomorphism of $\1$.

\begin{equation} \label{circlestd} \ig{1}{circlestd} \end{equation}

\begin{equation} \label{capcupstd} \ig{1}{capcupstd} \end{equation}

These two relations imply that cups and caps form inverse isomorphisms from $ss$ to $\1$. The following relation, which states that the $2m$-valent vertex gives an isomorphism from
$\ubr{sts\ldots}{m}$ to $\ubr{tst\ldots}{m}$, holds for any $s,t \in S$ with $m_{s,t} < \infty$.

\begin{equation} \label{2misomstd} \ig{1}{2misomstd} \end{equation}

The remaining relations come from triples $s,t,u \in S$ such that the corresponding parabolic subgroup is finite. By the classification of finite Coxeter groups the finite rank three parabolic subgroups can only be one of the following types:

\begin{eqnarray} (A_1 \times I_2(m)) && \ig{1}{A1I2mstd} \label{A1I2mstd} \end{eqnarray}

A specific example, when $m=2$, is:

\begin{eqnarray} (A_1 \times A_1 \times A_1) && \ig{1}{A1A1A1std} \label{A1A1A1std} \end{eqnarray}
	
\begin{eqnarray} (A_3) && \ig{1}{A3std} \label{A3std} \end{eqnarray}

\begin{eqnarray} (B_3) && \ig{1}{B3std} \label{B3std} \end{eqnarray}

\begin{eqnarray} (H_3) && \ig{1}{H3leftdashed} = \ig{1}{H3rightdashed} \label{H3std} \end{eqnarray}
This concludes the definition of $\DStd$. \end{defn}

\begin{defn} If we do not allow any grading shifts, and only allow $\Bbbk$-linear combinations of standard graphs without polynomials, we get a monoidal category $\DStd_0$ consisting entirely of degree $0$ maps. \end{defn}

\begin{remark} When working with $\DStd_0$, the base ring $\Bbbk$ is entirely arbitrary. When defining $\DStd$, as for $\StdBim$, one needs a realization of $W$ to obtain the ring $R$ but
no additional restrictions on $R$ are necessary. The main theorem \ref{FStdEquiv} below holds more generally for any (graded) ring $R$ with a faithful $W$-action.\end{remark}

Now we wish to show that $\DStd \cong \StdBim$, and that $\DStd_0 \cong \StdBim_0$.

\begin{defn} We define a functor $\FStd \co \DStd \to \StdBim$. Given an expression ${\un{x}} = st\dots u$, $\FStd$ sends the corresponding object to $R_s \ot R_t \ot \cdots \ot R_u$. The box containing $f$ is sent to multiplication by $f$. The cups, caps, and $2m$-valent vertices are all sent to the isomorphisms of $R$-bimodules which send 1-tensors to 1-tensors. This restricts to a functor
$\FStd_0 \co \DStd_0 \to \StdBim_0$. \end{defn}

\begin{claim} The functors $\FStd$ and $\FStd_0$ are well-defined. \end{claim}

\begin{proof} For $\FStd_0$, we need only check the relations minus (\ref{slidepolythrustd}). However, all relations hold immediately because both sides are bona fide maps in
$\StdBim_0$, and both sides send the 1-tensor to the 1-tensor. In addition, it is clear that equation (\ref{slidepolythrustd}) holds for standard bimodules, so $\FStd$ is well-defined.
\end{proof}

$\FStd$ is obviously essentially surjective. Suppose that the realization is faithful. Whenever a Hom space in $\StdBim$ is non-zero, it is generated by the isomorphism which sends the
1-tensor to the 1-tensor. This is clearly in the image of $\FStd$, so $\FStd$ is full.

\begin{thm} The functors $\FStd_0$ and $\FStd$ are equivalences, so long as the realization is faithful. \label{FStdEquiv} \end{thm}

This theorem is the main result of \cite{EWFenn}. Essentially, one needs to show that the space of morphisms from ${\un{x}}$ to ${\un{y}}$ is one-dimensional if $x=y$ (it is easy to see
that it is $0$-dimensional otherwise). Using biadjunction and various isomorphisms in both categories, this reduces to the fact that standard diagrams with empty boundary are polynomial
multiples of the empty diagram. The proof is actually topological in nature. Using Fenn's theory of diagrams \cite{Fenn} we are able to relate standard diagrams with empty boundary to
elements of the second homotopy group of a space related to the Coxeter complex.

We will discuss one important feature of the story in the next section.

%===========
\subsection{Rex moves and rex graphs}
\label{subsec-RexMoves}
%===========

The terminology of this section is ad hoc, and non-standard. This is due to ignorance, not malice.

For each $w \in W$, let $\G_w$ denote the set of all reduced expressions for $w$. This can be given the structure of a connected graph, the \emph{rex graph},  where two reduced expressions are connected by an edge if they differ by a single application of a braid relation \eqref{braidreln}. The edge itself can be labelled by the pair $\{s,t\}$ corresponding to the braid
relation; in other words, the edges are labelled by finitary rank 2 subsets of $S$. We shall only distinguish here between two different labellings: \emph{distant edges} for which
$m_{s,t}=2$, and those for which $m_{s,t}>2$. Let $\oG_w$ denote the graph obtained from $\G_w$ by contracting the distant edges.

We can associate a vertex ${\un{x}}$ of $\G_w$ with the standard bimodule $R_{\un{x}}$, isomorphic to $R_w$. We think of a path as a \emph{rex move}, a sequence of braid relations traversing the rex graph to a new reduced expression. To an
edge labelled $\{s,t\}$ from ${\un{x}}$ to ${\un{y}}$, we can associate the $2m_{s,t}$-valent vertex, which is a morphism from $R_{\un{x}}$ to $R_{\un{y}}$. Therefore we can associate a morphism in $\DStd_0$
to each path in $\G_w$, and we also call this morphism a rex move. We can never construct a morphism with caps or cups in this fashion, since those involve non-reduced expressions. In
fact, $\DStd$ essentially encodes a study of the graph of all expressions, reduced and non-reduced, for $w$; this point is made clear in \cite{EWFenn}, and may become more obvious after
the remainder of this section.

The relations of $\DStd$ which do not involve cups and caps, namely \eqref{2misomstd} through \eqref{H3std}, all come from loops (i.e. paths from a vertex to itself) in $\G_w$. That is,
each side of the relation comes from a path from $\un{x}$ to $\un{x}'$ in $\G_w$, and the relation states that the corresponding morphisms are equal. Equivalently, once \eqref{2misomstd}
is known, the relation states that the loop $\un{x} \to \un{x}' \to \un{x}$ is equal to the trivial loop at $\un{x}$.

For each edge there is a \emph{boring} loop which follows the edge and then follows it in reverse, which corresponds to \eqref{2misomstd}. In addition, each finitary rank 3 subset of $S$
gives rise to a kind of loop, which first appears in the longest elements of these parabolic subgroups.

\begin{ex} We give two examples of type $A_1 \times I_2(m)$. The labeling of indices in $S$ should be obvious. Type $A_1 \times A_1 \times A_1$: \igc{1}{135graph} Type $A_1 \times
A_2$: \igc{1}{1214graph} \end{ex}

\begin{ex} Here is $\G_{w_0}$ for type $A_3$. The orientations on the arrows will be explained at the end of this section. \igc{1}{121321graph} \end{ex}

\begin{ex} Here is $\oG_{w_0}$ for type $B_3$. This can be deduced from \eqref{B3std}, and we let the avid reader do the same for $H_3$. The red and blue edges correspond to different
parabolic subsets. \igc{1}{B3rexgraph} \end{ex}

There is one additional kind of loop, which arises when two braid relations can be applied to disjoint parts of a rex. This is called a \emph{disjoint square}. In $\DStd$, this corresponds
to the fact that distant pictures commute in a diagrammatic category. Disjoint squares need not involve disjoint colors.

\begin{ex} A disjoint square.  \igc{1}{disjointsquare} \end{ex}

It is known \cite[Chapter 2, \S 5]{Ronan} that these loops generate the set of all loops in $\G_w$, in a sense which the reader can intuit. This corresponds topologically to the fact that
$\pi_2$ of the truncated dual Coxeter complex is generated by the finitary rank 3 subsets. In other words, the relations \eqref{2misomstd} through \eqref{H3std} are sufficient to imply
that any two paths from $\un{x}$ to $\un{x}'$ will induce the same morphism in $\DStd$. The theory of reduced expressions is enough to say that diagrams in $\DStd$ without cups and caps form a one-dimensional space. The proof in \cite{EWFenn} deals with non-reduced expressions as well.

Let us prepare the reader for the following chapter. Note that the only loops in $\oG_w$ (in addition to the boring loops and disjoint squares) come from \emph{connected} finitary rank 3
subsets of $S$, which are of the type $A_3$, $B_3$, or $H_3$. The equality of the two sides in \eqref{A3std} is often called the \emph{Zamolodzhikov tetrahedron equation}, and so we refer to the
$B3$ and $H3$ equations as Zamolodzhikov equations as well.

The story will be complicated much further when we start working with Bott-Samelson bimodules instead of standard bimodules. For two rexes ${\un{x}}$ and ${\un{y}}$ connected by an edge
labelled by the subset $J=\{s,t\}$, there will be a corresponding map from $B_{\un{x}}$ to $B_{\un{y}}$ which projects from $\ubr{B_s \ot B_t \ot \cdots}{m} \to B_{w_J}$ and then includes
$B_{w_J} \to \ubr{B_t \ot B_s \ot \cdots}{m}$. To any path in $\G_w$ we can still associate a morphism in $\BSBim$, but two paths will not always agree. There is no \eqref{2misomstd}; the
boring loop is projection to a summand, not the identity map. It is true, at least, that following the boring loop twice is the same as following it once. After localizing, however, rex
moves in $\BSBim$ restricted to the unique ``top" summand $Q_w$ will give exactly the rex moves in $\DStd$. In particular, this implies that two paths give rex moves in $\BSBim$ which
agree modulo ``lower terms", where lower terms are maps which, when localized, give the zero map on $Q_w$. For any ${\un{x}},{\un{y}} \in \G_w$, the space of maps from $B_{\un{x}}$ to
$B_{\un{y}}$ modulo lower terms is free of rank 1 over $R$, as we shall see.

% \begin{remark}
% Manin and Schechtmann \cite{MS} have a beautiful study of the graphs $\G_w$ and $\oG_w$ in type $A$. They showed that the graph $\oG_w$ can be oriented (lexicographically), in such a
% way that the graph has a unique source and a unique sink, and any two oriented paths between two vertices have the same length. Any two oriented paths are related by the oriented $A_3$ Zamalodchikov relation, or by disjoint squares.
% 
% In fact, the story continues! The set of oriented paths from $\un{x}$ to $\un{y}$ forms a graph, with two kinds of edges corresponding to disjoint squares and to $A_3$ loops. Contracting the disjoint edges, we get a graph of ``paths between paths'' which can be oriented with a unique source and sink. Two oriented paths in this graph are related by either a disjoint square or a loop corresponding to $A_4$. The set of oriented paths in this graph can be dealt with the same way, with loops coming from $A_5$, and so forth. The original graph $\G_w$ was itself the set of oriented paths from the identity to $w$ in the graph of all elements of $W$, and the original edges came from a loop associated to $A_2$.
% 
% It is an amazing fact that the two \emph{oriented} paths in the $A_3$ loop agree as morphisms in $\BSBim$. If the three vertices in $A_3$ are $J \subset S$, then the two longest
% oriented paths are a sequence of projections to and inclusions of summands, which overall manage to express the projection-inclusion $B_{\un{x}} \to B_J \to B_{\un{y}}$. However, two non-oriented
% paths which together form the $A_3$ loop need \emph{not} agree in $\BSBim$. For more discussion of this, and of \cite{MS} in general, see \cite{EInd}.
% 
% Reduced expressions in type $B$ admit similar lexicographic orientations, which permit unique sources and sinks. Extending this lexicographic order to ``paths between paths" in a similar
% fashion is currently under investigation. The previous paragraph applies to the $B_3$ loop precisely as it does to the $A_3$ loop. However, we do not believe that rex graphs in $H_3$
% should admit a nice orientation. There is (alas!) no preferred source ${\un{x}}$ in the rex graph for the longest element of $H_3$ such that the two paths from ${\un{x}}$ to its opposite
% point ${\un{y}}$ give equal rex moves in $\BSBim$. \end{remark}